% Circuito completo del TP2 (3 x TL084, alimentacion +/-12 V).
% Se incluye con \input dentro de un figure*. Coordenadas absolutas.
\begin{circuitikz}[scale=0.62, transform shape, font=\small]
\ctikzset{resistors/scale=0.7, capacitors/scale=0.7, diodes/scale=0.7,
          tripoles/op amp/width=1.5, tripoles/op amp/height=1.5,
          tripoles/op amp/input height=0.5, bipoles/length=1.2cm}
\tikzset{et/.style={font=\footnotesize\bfseries, color=blue!55!black},
         senal/.style={font=\normalsize\bfseries, color=red!70!black},
         ref/.style={font=\small}}

% =============================== FILA 1 ===============================
% ---- Entrada y adaptacion de impedancia (U1A) ----
\node[et, align=center] at (2.4,3.3) {Adaptación de $Z_i$};
\draw (0,0) node[ocirc] {} node[above, align=center, font=\footnotesize, yshift=2pt]
      {Generador\\0 a \SI{2,5}{\volt}\\\SI{4300}{\hertz}};
\draw (0.08,0) to[short, -*] (0.9,0) to[R, l_={$R_1$\\\SI{15}{\kilo\ohm}}, align=left] (0.9,-1.8) node[ground] {};
\node[op amp, noinv input up, anchor=+] (u1a) at (2.0,0) {\scriptsize U1A};
\draw (0.9,0) -- (u1a.+);
\draw (u1a.-) -- (1.7,-1.0) -- (1.7,-1.7) -- (4.4,-1.7) -- (4.4,-0.5);
\draw (u1a.out) to[short, -*] (4.4,-0.5);

% ---- Senal 1: amplificador diferencial (U1C) ----
\node[et] at (8.0,3.3) {Señal 1};
\draw (4.4,-0.5) to[R, l={$R_2$ \SI{10}{\kilo\ohm}}] (6.4,-0.5) to[short, -*] (6.7,-0.5);
\node[op amp, anchor=-] (u1c) at (7.2,0.5) {\scriptsize U1C};
\draw (6.7,-0.5) -- (u1c.+);
\draw (6.7,-0.5) to[R, l_={$R_3$\\\SI{10}{\kilo\ohm}}, align=right] (6.7,-2.3) node[ground] {};
\draw (u1c.-) to[short, -*] (6.9,0.5);
\draw (6.9,0.5) to[R, l_={$R_4$ \SI{10}{\kilo\ohm}}] (5.0,0.5) node[left, ref] {$V_{\mathrm{REF1}}$};
\draw (6.9,0.5) -- (6.9,1.6) to[R, l={$R_5$ \SI{10}{\kilo\ohm}}] (9.6,1.6) -- (9.6,0) -- (u1c.out);
\draw (9.6,0) to[short, *-*] (10.0,0);
\node[senal, below] at (10.0,-0.1) {S1};

% ---- Integrador con acoplamiento capacitivo (U1D) ----
\node[et] at (14.6,3.3) {Integrador};
\draw (10.0,0) to[C, l={$C_C$ \SI{1}{\micro\farad}}] (11.7,0)
      to[R, l={$R_I$ \SI{5,6}{\kilo\ohm}}] (13.4,0) to[short, -*] (13.7,0);
\node[op amp, anchor=-] (u1d) at (14.2,0) {\scriptsize U1D};
\draw (13.7,0) -- (u1d.-);
\draw (13.7,0) -- (13.7,1.1) to[C, l={$C_1$ \SI{4,7}{\nano\farad}}] (16.6,1.1);
\draw (13.7,1.1) to[short, *-] (13.7,2.2) to[R, l={$R_F$ \SI{1}{\mega\ohm}}] (16.6,2.2) -- (16.6,1.1) to[short, *-] (16.6,-0.5) -- (u1d.out);
\draw (u1d.+) -- (13.9,-1.0) node[left, ref] {$V_{\mathrm{REF2}}$};
\draw (16.6,-0.5) to[short, *-] (16.9,-0.5);

% ---- Senal 2: inversor alrededor de VREF2 (U2A) ----
\node[et] at (20.6,3.3) {Señal 2};
\draw (16.9,-0.5) to[R, l={$R_6$ \SI{10}{\kilo\ohm}}] (18.7,-0.5) to[short, -*] (19.0,-0.5);
\node[op amp, anchor=-] (u2a) at (19.5,-0.5) {\scriptsize U2A};
\draw (19.0,-0.5) -- (u2a.-);
\draw (19.0,-0.5) -- (19.0,0.6) to[R, l={$R_7$ \SI{10}{\kilo\ohm}}] (21.9,0.6) -- (21.9,-1.0) -- (u2a.out);
\draw (u2a.+) -- (19.2,-1.5) node[left, ref] {$V_{\mathrm{REF2}}$};
\draw (21.9,-1.0) to[short, *-*] (22.4,-1.0);
\node[senal, right] at (22.5,-1.0) {S2};

% =============================== FILA 2 ===============================
% ---- Referencia de precision con buffer (DZ1, P1, P2, U1B, U2D) ----
\node[et, align=center] at (3.0,-3.0) {Referencia de precisión};
\draw (0.6,-3.7) node[vcc] {$+\SI{12}{\volt}$} to[R, l_={$R_Z$\\\SI{470}{\ohm}}, align=left] (0.6,-5.2);
\draw (0.6,-7.0) node[ground] {} to[zDo, l^={$D_{Z1}$\\1N750\\\SI{4,7}{\volt}}, align=right, -*] (0.6,-5.2);
\draw (0.6,-5.2) to[short, -*] (1.3,-5.2) -- (1.9,-5.2) to[pR, n=pa, l_={$P_1$}, mirror] (1.9,-7.0) node[ground] {};
\draw (1.3,-5.2) -- (1.3,-8.3) -- (1.9,-8.3) to[pR, n=pb, l_={$P_2$}, mirror] (1.9,-10.1) node[ground] {};
\node[op amp, noinv input up, anchor=+] (u1b) at (3.0,-6.1) {\scriptsize U1B};
\draw (pa.wiper) -| (2.7,-6.1) -- (u1b.+);
\draw (u1b.-) -- (2.75,-7.1) -- (2.75,-7.6) -- (5.4,-7.6) -- (5.4,-6.6) -- (u1b.out);
\draw (5.4,-6.6) to[short, *-o] (5.8,-6.6) node[right, ref] {$V_{\mathrm{REF2}}$};
\node[op amp, noinv input up, anchor=+] (u2d) at (3.0,-9.2) {\scriptsize U2D};
\draw (pb.wiper) -| (2.7,-9.2) -- (u2d.+);
\draw (u2d.-) -- (2.75,-10.2) -- (2.75,-10.7) -- (5.4,-10.7) -- (5.4,-9.7) -- (u2d.out);
\draw (5.4,-9.7) to[short, *-o] (5.8,-9.7) node[right, ref] {$V_{\mathrm{REF1}}$};
\node[font=\footnotesize, align=left, anchor=west] at (0.0,-11.6)
     {$P_1$, $P_2$: preset multivuelta \SI{10}{\kilo\ohm}\\
      $V_{\mathrm{REF2}} = \SI{1,25}{\volt}$, $V_{\mathrm{REF1}} = \SI{1,00}{\volt}$};

% ---- Derivador compensado (U2B) ----
\node[et] at (10.3,-3.5) {Derivador};
\draw (22.4,-1.0) -- (22.4,-2.8) -- (9.0,-2.8) -- (9.0,-5.8);
\draw (9.0,-5.8) to[R, l_={$R_S$ \SI{1}{\kilo\ohm}}] (10.6,-5.8)
      to[C, l_={$C_D$ \SI{1}{\nano\farad}}] (12.2,-5.8) to[short, -*] (12.5,-5.8);
\node[op amp, anchor=-] (u2b) at (13.0,-5.8) {\scriptsize U2B};
\draw (12.5,-5.8) -- (u2b.-);
\draw (u2b.+) -- (12.7,-6.8) node[ground] {};
\draw (12.5,-5.8) -- (12.5,-4.8) to[C, l={$C_F$ \SI{56}{\pico\farad}}] (15.4,-4.8);
\draw (12.5,-4.8) to[short, *-] (12.5,-3.8) to[R, l={$R_D$ \SI{31,6}{\kilo\ohm}}] (15.4,-3.8) -- (15.4,-4.8) to[short, *-] (15.4,-6.3) -- (u2b.out);
\draw (15.4,-6.3) to[short, *-] (15.7,-6.3);

% ---- Senal 3: inversor con desplazamiento (U2C) ----
\node[et] at (19.4,-3.5) {Señal 3};
\draw (15.7,-6.3) to[R, l={$R_{10}$ \SI{10}{\kilo\ohm}}] (17.5,-6.3) to[short, -*] (17.8,-6.3);
\node[op amp, anchor=-] (u2c) at (18.3,-6.3) {\scriptsize U2C};
\draw (17.8,-6.3) -- (u2c.-);
\draw (17.8,-6.3) -- (17.8,-5.2) to[R, l={$R_{11}$ \SI{10}{\kilo\ohm}}] (20.7,-5.2) -- (20.7,-6.8) -- (u2c.out);
\draw (u2c.+) to[short, -*] (17.9,-7.3);
\draw (17.9,-7.3) to[R, l_={$R_9$\\\SI{15}{\kilo\ohm}}, align=right] (17.9,-9.0) node[ground] {};
\draw (17.9,-7.3) to[R, l_={$R_8$ \SI{10}{\kilo\ohm}}] (16.0,-7.3) node[left, ref] {$V_{\mathrm{REF2}}$};
\draw (20.7,-6.8) to[short, *-*] (21.1,-6.8);
\node[senal, right] at (21.2,-6.8) {S3};

% =============================== FILA 3 ===============================
% ---- Generador de impulsos: pasaaltos activo (U3A) ----
\node[et, align=center] at (13.6,-14.9) {Generador de impulsos};
\draw (21.1,-6.8) -- (21.1,-9.8) -- (8.6,-9.8) -- (8.6,-12.5);
\draw (8.6,-12.5) to[R, l_={$R_{S4}$ \SI{4,7}{\kilo\ohm}}] (10.3,-12.5)
      to[C, l_={$C_4$ \SI{1}{\nano\farad}}] (11.9,-12.5) to[short, -*] (12.2,-12.5);
\node[op amp, anchor=-] (u3a) at (12.7,-12.5) {\scriptsize U3A};
\draw (12.2,-12.5) -- (u3a.-);
\draw (u3a.+) -- (12.4,-13.5) node[ground] {};
\draw (12.2,-12.5) -- (12.2,-11.4) to[R, l={$R_{F4}$ \SI{4,7}{\kilo\ohm}}] (15.1,-11.4) -- (15.1,-13.0) -- (u3a.out);
\draw (15.1,-13.0) to[short, *-*] (15.5,-13.0);
\node[font=\small, below left] at (15.5,-13.0) {$v_p$};

% ---- Rectificador de precision (U3B) ----
\node[et, align=center] at (19.4,-16.4) {Rectificador de precisión};
\draw (15.5,-13.0) to[R, l_={$R_{12}$ \SI{10}{\kilo\ohm}}] (17.3,-13.0) to[short, -*] (17.6,-13.0);
\node[op amp, anchor=-] (u3b) at (18.2,-13.0) {\scriptsize U3B};
\draw (17.6,-13.0) -- (u3b.-);
\draw (u3b.+) node[ground] {};
\draw (u3b.out) to[short, -*] (20.6,-13.5);
\draw (20.6,-13.5) -- (20.6,-11.8) to[D, l^={$D_1$ 1N4004}] (17.6,-11.8) -- (17.6,-13.0);
\draw (17.6,-13.0) -- (17.6,-15.2) to[R, l_={$R_{13}$ \SI{10}{\kilo\ohm}}] (21.4,-15.2);
\draw (21.4,-15.2) to[D, l_={$D_2$\\1N4004}, align=left, *-] (21.4,-13.5) -- (20.6,-13.5);
\node[font=\small, below left] at (21.4,-15.2) {$v_A$};

% ---- Sumador de salida (U3C): senal 4 ----
\node[et] at (25.9,-10.9) {Sumador};
\draw (15.5,-13.0) -- (15.5,-10.4) to[R, l_={$R_{14}$ \SI{20}{\kilo\ohm}}] (23.4,-10.4) -- (23.4,-13.0);
\draw (21.4,-15.2) to[R, l_={$R_{15}$ \SI{10}{\kilo\ohm}}] (23.4,-15.2) -- (23.4,-13.0);
\node[op amp, noinv input up, anchor=-] (u3c) at (24.2,-13.0) {\scriptsize U3C};
\draw (23.4,-13.0) to[short, *-*] (23.8,-13.0) -- (u3c.-);
\draw (u3c.+) -- (23.9,-12.0) node[ground] {};
\draw (23.8,-13.0) -- (23.8,-16.4) to[R, l_={$R_{F8}$ \SI{10}{\kilo\ohm}}] (25.3,-16.4)
      to[vR, l_={$P_3$ \SI{10}{\kilo\ohm}}] (26.8,-16.4) -- (26.8,-12.5) -- (u3c.out);
\draw (26.8,-12.5) to[short, *-o] (27.3,-12.5);
\node[senal, right] at (27.4,-12.5) {S4};

% ---- Notas ----
\node[font=\footnotesize, align=left, anchor=north west] at (0.0,-13.2)
     {U1, U2, U3: TL084 (U3D sin uso).\\
      Alimentación: $\pm\SI{12}{\volt}$, desacople de\\
      \SI{100}{\nano\farad} + \SI{10}{\micro\farad} en cada integrado.\\
      Resistores de película metálica, \SI{1}{\percent}.};
\end{circuitikz}
